\documentclass{article}

% if you need to pass options to natbib, use, e.g.:
%     \PassOptionsToPackage{numbers, compress}{natbib}
% before loading neurips_2019

% ready for submission
% \usepackage{neurips_2019}

% to compile a preprint version, e.g., for submission to arXiv, add add the
% [preprint] option:
%     \usepackage[preprint]{neurips_2019}

% to compile a camera-ready version, add the [final] option, e.g.:
\usepackage[final]{neurips_2019}

% to avoid loading the natbib package, add option nonatbib:
%     \usepackage[nonatbib]{neurips_2019}

\usepackage[utf8]{inputenc} % allow utf-8 input
\usepackage[T1]{fontenc}    % use 8-bit T1 fonts
\usepackage{hyperref}       % hyperlinks
\usepackage{url}            % simple URL typesetting
\usepackage{booktabs}       % professional-quality tables
\usepackage{amsfonts}       % blackboard math symbols
\usepackage{nicefrac}       % compact symbols for 1/2, etc.
\usepackage{microtype}      % microtypography
\usepackage{graphicx}

\title{Final report template}

% The \author macro works with any number of authors. There are two commands
% used to separate the names and addresses of multiple authors: \And and \AND.
%
% Using \And between authors leaves it to LaTeX to determine where to break the
% lines. Using \AND forces a line break at that point. So, if LaTeX puts 3 of 4
% authors names on the first line, and the last on the second line, try using
% \AND instead of \And before the third author name.

\author{%
  Author 1\\
  Department of Computer Engineering\\
  Iran University of Science \\ and Technology\\
  \texttt{name@cs.iust.ac.ir} \\
  % examples of more authors
  % \And
  % Coauthor \\
  % Affiliation \\
  % \texttt{email} \\
  % \AND
  % Coauthor \\
  % Affiliation \\
  % Address \\
  % \texttt{email} \\
}

\begin{document}


\noindent\begin{minipage}{0.1\textwidth}% adapt widths of minipages to your needs
\includegraphics[width=1.1cm]{iust_logo.png}
\end{minipage}%
\hfill%
\begin{minipage}{1\textwidth}\raggedright
Iran University of Science and Technology\\
Deep Learning (Spring 2019)\\
\end{minipage}
% \end{}

\maketitle

\begin{abstract}
  You should give a brief introduction of your problem and the solution that you have provided.
\end{abstract}

\section{Introduction}
 In this part, you should explain the problem or the task you want to solve and demonstrate the challenges that you may face during solving this problem.

\section{Related work/Background}
A concise explanation of two to three recent works that had proposed an approach to overcome or solve the problem. These works might have found a solution for different settings of this problem.

\section{Proposed method}
In this section, you must provide a comprehensive explanation of your final improved method to solve this problem. Some of the things that you should describe in the report are the following items:

\begin{itemize}
    \item Any kind of preprocessing or normalization (if you have used).
    \item The architecture of the neural network that you have used. It is better to demonstrate this with a graph.
\end{itemize}{}

\section{Results}
At first, you should describe the dataset that you have used, and then, you must describe the experiments that you maid, and its results. It is better to test your proposed method in different settings and with various parameters.

\section{Discussion}
In this section, you will discuss your proposed method to solve the problem. Some of the questions that you might find useful to answer are:
\begin{itemize}
    \item Why is this specific type of Deep Neural Networks suitable for this task?
    \item What are the benefits of your proposed model against other approaches?
\end{itemize}

\section*{References}

\medskip

\small

[1] Alexander, J.A.\ \& Mozer, M.C.\ (1995) Template-based algorithms for
connectionist rule extraction. In G.\ Tesauro, D.S.\ Touretzky and T.K.\ Leen
(eds.), {\it Advances in Neural Information Processing Systems 7},
pp.\ 609--616. Cambridge, MA: MIT Press.

[2] Bower, J.M.\ \& Beeman, D.\ (1995) {\it The Book of GENESIS: Exploring
  Realistic Neural Models with the GEneral NEural SImulation System.}  New York:
TELOS/Springer--Verlag.

[3] Hasselmo, M.E., Schnell, E.\ \& Barkai, E.\ (1995) Dynamics of learning and
recall at excitatory recurrent synapses and cholinergic modulation in rat
hippocampal region CA3. {\it Journal of Neuroscience} {\bf 15}(7):5249-5262.

\end{document}
